\chapter{Machine Learning for Trading}
\label{ch:ml_trading}
\chaptermark{ML for trading}

Most ML-for-trading projects lose money because they treat the
problem as ``forecast tomorrow's return with a neural net'':
predicting a near-random-walk target with a high-variance estimator,
then being surprised when out-of-sample performance is indistinguishable
from noise. The real opportunities are narrower: feature engineering on
order-book data, regime detection, trade-level \emph{meta-labeling}, and
execution. This chapter draws the bulk of its content from Lop\'ez de
Prado's \textit{Advances in Financial Machine Learning} (AFML), a book
less about ``how to build trading ML models'' than ``here is every way
people get it wrong.''

\Cref{ch:backtesting} spent thirty pages on validation discipline, and
every sentence of this chapter assumes you know why naive $k$-fold
cross-validation on financial time series is a bug and what CPCV does
about it. The techniques here produce nonsense if fed into a broken
validation pipeline.

\begin{backgroundbox}[Prerequisites]
From earlier in the book:
\begin{itemize}
  \item \Cref{ch:mean_reversion_ou}: Ornstein--Uhlenbeck mean
  reversion. The ``OU $z$-score entry'' primary signal is the
  running example for meta-labeling in
  \cref{sec:ml_metalabel}.
  \item \Cref{ch:pairs_cointegration}: pairs-trading mechanics,
  used as the worked example for the meta-labeling lift
  calculation.
  \item \Cref{ch:backtesting}: Combinatorial Purged
  Cross-Validation (CPCV), walk-forward, and the general
  pitfalls of look-ahead and data snooping. This chapter
  refers to CPCV as the correct inner validation loop for every
  ML procedure it introduces.
  \item \Cref{ch:kelly_risk}: Sharpe, drawdown, and the mental
  model of ``scoring after sizing''. Needed for the
  meta-labeling lift numbers in \cref{sec:ml_metalabel}.
\end{itemize}
From outside the book:
\begin{itemize}
  \item Elementary supervised learning: linear/logistic
  regression, decision trees, random forests, basic
  classification metrics (precision, recall, ROC-AUC), and the
  idea of cross-validation for hyperparameter selection.
  \item The concept that \emph{financial time series are not
  stationary}: distributions drift, regimes change, and the
  i.i.d.\ assumption powering textbook ML is always false for
  markets. This chapter is about what to do given that fact.
\end{itemize}
\end{backgroundbox}

% =====================================================================
\section{Where ML helps}
\label{sec:ml_where_helps}

ML does not help with ``what will the price of product $X$ be
tomorrow.'' It helps with four narrower questions, and in each case
the reason it helps is not the cleverness of the model but the
structure of the prediction target. We take them in order of
practical impact.

\subsection{Feature engineering on order-book data}
\label{subsec:ml_feature_eng}

The most productive use of ML in trading is \emph{compression},
not prediction. A Level-2 order book at millisecond frequency is
a firehose of numbers: ten levels per side, quantity and price at
each, plus every trade print. A rules-based strategy can use at most
a handful of features before it becomes unmaintainable. ML reduces
hundreds of raw order-book columns into a small number of short-horizon
predictive features.

Feeding the raw order-book columns directly into a classifier is a
non-starter: raw prices drift (non-stationary), quantity columns
change variance with the regime (heteroscedastic), and both swamp any
learner with scale artefacts rather than signal. Good features are
\emph{stationary transforms} (log-returns, rolling $z$-scores,
normalised imbalances) whose distribution is comparable across
time and across instruments.

Examples of such derived features, all covered in
\cref{ch:order_flow}:

\begin{itemize}
  \item \textbf{Order Flow Imbalance (OFI)}: signed running total
  of aggressive buy minus aggressive sell volume over a
  lookback window. Empirically one of the strongest
  short-horizon return predictors.
  \item \textbf{Top-of-book imbalance}:
  $(q^{\mathrm{bid}}_1 - q^{\mathrm{ask}}_1) / (q^{\mathrm{bid}}_1 + q^{\mathrm{ask}}_1)$
  at the inside quote. A cheap real-time drift proxy.
  \item \textbf{Realised volatility at multiple horizons}:
  $\volat_{1\mathrm{m}}, \volat_{5\mathrm{m}},
  \volat_{30\mathrm{m}}$ from intraday returns. Used as direct
  features and to normalise other features across regimes.
  \item \textbf{Trade direction entropy}: Shannon entropy of
  signed trades over a rolling window. Low entropy = one-sided
  flow; high = balanced flow.
  \item \textbf{Queue-position-weighted imbalance}: imbalance
  that discounts orders at the back of a price-time priority
  queue, which are less likely to fill soon.
\end{itemize}

The ML part is boring on purpose. A gradient-boosted tree or linear
model predicts the one-step-ahead return from these features under
CPCV. The value added is not the model (a linear regression
usually does fine) but the \emph{choice of features}. A skilled
researcher spends 80\% of their time on feature construction and
20\% on the model. Academic papers on ``deep learning for financial
prediction'' invert that ratio and are rarely reproducible.

\subsection{Regime detection}
\label{subsec:ml_regime}

The second win is regime detection: given a history of market
observables (volatility, spread, volume, correlation structure),
infer how many distinct regimes the market has been in and which it
is in now. The canonical tool is the \textbf{Hidden Markov Model}
(HMM), which assumes the market cycles between a small number of
latent states (``quiet mean-reverting,'' ``high-vol trending,''
``crisis'') with observable features drawn from a state-dependent
distribution. \Cref{ch:which_strategy_when} promotes regime detection
from ``nice to have'' to the backbone of the ``which strategy when''
framework.

Regime detection is a good ML use case because the target is not a
scalar future return but a discrete label that can be validated
qualitatively (``yes, March 2020 was a crisis regime''). The
prediction problem is more stable than return prediction: volatility
is far more auto-correlated than the return itself, so an HMM
fitted on year $t$ still works on year $t+1$ provided the training
regimes also appear at test time.

\subsection{Trade-level meta-labeling}
\label{subsec:ml_metalabel_intro}

The third, and in Lop\'ez de Prado's framing the highest-value
ML use in trading, is \textbf{meta-labeling}: given a primary
strategy that emits trade signals, train a binary classifier to
decide \emph{which signals to actually take}. The primary strategy
supplies direction (long/short, entry/exit); the meta-labeler
supplies the filter. We treat this as its own section
(\cref{sec:ml_metalabel}) because its payoff relative to
implementation cost is unmatched.

\subsection{Execution}
\label{subsec:ml_execution}

The fourth win is execution: reinforcement learning (RL) for
optimal order placement and routing. Given a parent order, an RL
agent learns a slicing and child-order policy that minimises
implementation shortfall relative to the arrival price. The
Almgren--Chriss solution from \cref{ch:almgren_chriss} is the
analytical baseline; RL relaxes the linear-impact assumption and
learns a policy respecting the empirically measured cost curve and
current order-book state. Execution is the area where RL papers
translate into genuine production lift, because the reward signal
is high-frequency, unambiguous, and approximately stationary over
short horizons.

We do not cover execution ML here (not on the Prosperity critical
path); interested readers should start with
\citet{cartea2015algorithmic} chapters 6--8 and the RL literature
that relaxes Almgren--Chriss's linear-impact simplification.

% =====================================================================
\section{Where ML fails}
\label{sec:ml_where_fails}

Having established the four regions where ML earns its keep, we
enumerate the larger regions where it does not. These are
categorical failure modes where more compute and a bigger model
actively make things worse.

\subsection{Naked return prediction}
\label{subsec:ml_naked_return}

The archetypal failed project is ``predict $r_{t+1}$ from
$(r_t, r_{t-1}, \ldots, r_{t-L})$ using a neural network.'' Two
reasons this fails, and they compound.

First, the signal-to-noise ratio is catastrophically low. A
one-day equity return has annualised volatility around $20\%$, so
a single daily return has standard deviation about $1.26\%$. A
real alpha signal generating Sharpe $1$ per year contributes a
daily drift of $\approx 0.08\%$: noise is $16\times$ the signal
on any single observation. Neural networks with millions of
parameters are wonderful at memorising the noise.

Second, one-period equity returns are a near-random-walk target
by construction: if they were not, a trivial linear predictor
would exist and be arbitraged away. The weak form of the
Efficient Market Hypothesis says precisely that an ML model
trained on lagged returns should find nothing; the fact that it
does find something in-sample is a diagnostic of overfitting.

The diagnosis is nearly always the same: stunning in-sample
$R^2$, a decent CV $R^2$ under the wrong CV scheme, and a
zero-or-negative $R^2$ on a genuinely held-out period. ``The
wrong CV scheme'' is almost always naive $k$-fold, and fixing it
kills the result.

\begin{pitfallbox}[Don't use $k$-fold CV on time series]
This is the most common ML-for-trading bug, responsible for a
huge fraction of irreproducible papers. The failure mode:

\begin{enumerate}
  \item You build triple-barrier labels
  (\cref{sec:ml_triple_barrier}) with a $5$-day horizon: the
  label at time $t$ depends on prices between $t$ and $t+5$.
  \item You apply scikit-learn's \texttt{KFold} with $k = 5$,
  which shuffles observations and splits them
  contiguously (or worse, randomly).
  \item A training observation at time $t' = t + 3$ has a label
  whose window overlaps with the test observation at time
  $t$. Information about \emph{the test label} has leaked
  into the training set via the training labels.
  \item The classifier reports a miraculous out-of-sample
  accuracy, e.g.\ $75\%$ on a 55\%-baseline problem.
  \item In live trading the accuracy collapses to the true
  55\% and the strategy loses money.
\end{enumerate}

Fix: CPCV (\cref{ch:backtesting}). Purge training observations
whose label window intersects a test observation, plus an embargo
on each side of the test fold. Any non-CPCV-compliant ML-trading
evaluation is, until proven otherwise, a leakage artefact.
\end{pitfallbox}

\subsection{Hyperparameter abuse}
\label{subsec:ml_hyper_abuse}

The second failure mode is hyperparameter abuse: tuning twenty
hyperparameters on the same data used for evaluation, the
data-snooping bias of \cref{ch:backtesting} in ML clothing. If you
run $N = 1000$ parameter combinations and report the best CV
Sharpe, the expected-max-of-$N$-iid-normals approximation
$\E[\max_{i \le N} Z_i] \approx \sqrt{2 \ln N}$ gives a spurious
Sharpe of $\sqrt{2 \ln 1000} \approx 3.72$ standard errors above
the true mean \emph{even if every strategy is pure noise}. The
defence is nested CPCV: outer CPCV for evaluation, a separate
inner CPCV inside each outer training fold for tuning.

Most applied papers skip the outer loop and report the best
inner-CV score as ``out-of-sample performance.'' This is silently
but devastatingly wrong.

\subsection{Non-stationarity}
\label{subsec:ml_nonstationary}

The third failure mode is baked into the data. Return
distributions drift: mean, variance, correlation structure, tail
behaviour, and the \emph{predictability structure} of features
all change as regimes change. A model with Sharpe $2$ on
$2015$--$2018$ may not survive $2020$ because the regime it
learned no longer exists. This is structural, not a bias to
correct.

Defences: (i)~retrain frequently on a rolling window, (ii)~pool
across instruments for statistical power, (iii)~condition on a
regime label from a separate detector (\cref{subsec:ml_regime}).
None fully solve the problem; they make degradation more gradual.
Expecting a fit-once-and-forget model to work for years is the
common unforced error.

\subsection{The general taxonomy}
\label{subsec:ml_general_taxonomy}

\begin{keyfactbox}[Where ML for trading lives and dies]
Four regions where ML helps:
\begin{enumerate}
  \item Feature engineering on order-book data
  (\cref{ch:order_flow}).
  \item Regime detection via HMMs and state-space models
  (\cref{ch:which_strategy_when}).
  \item Trade-level meta-labeling of a rule-based primary
  strategy (\cref{sec:ml_metalabel}).
  \item Execution: RL for order slicing and routing
  (\cref{ch:almgren_chriss} as baseline).
\end{enumerate}
Four regions where ML fails:
\begin{enumerate}
  \item Naked return prediction from lagged returns.
  \item $k$-fold CV on time series (fixed by CPCV).
  \item Hyperparameter abuse (fixed by nested CPCV).
  \item Non-stationarity (mitigated, never eliminated, by
  rolling retraining and regime conditioning).
\end{enumerate}
\end{keyfactbox}

% =====================================================================
\section{Triple-barrier labeling}
\label{sec:ml_triple_barrier}

Before training any classifier we need labels. The textbook
approach (``$+1$ if the next-step return is positive, $-1$
otherwise'') is wrong for trading. First, it assumes every
trade is held for exactly one step. Second, it ignores risk: a
$0.01\%$ return is labelled identically to a $5\%$ return. The
\textbf{triple-barrier method} of \citeAFML{} Ch.~3 fixes both.

\subsection{The method}
\label{subsec:ml_tb_method}

At each observation time $t$, suppose the current price is
$p_t$ and the estimated short-horizon volatility is $\hat\volat_t$
(e.g.\ an exponentially weighted standard deviation of daily
returns). We open a hypothetical trade at $p_t$ and set three
barriers:

\begin{itemize}
  \item \textbf{Top barrier} (profit target):
  $p^{\mathrm{top}}_t = p_t (1 + \pi_u \hat\volat_t)$, for some
  profit-target multiplier $\pi_u > 0$.
  \item \textbf{Bottom barrier} (stop-loss):
  $p^{\mathrm{bot}}_t = p_t (1 - \pi_d \hat\volat_t)$, for some
  stop-loss multiplier $\pi_d > 0$.
  \item \textbf{Vertical barrier} (max holding): we give up on
  the trade after $H$ time steps and record the realised return
  regardless.
\end{itemize}

The label $y_t$ is then the first barrier the price path hits:

\[
y_t = \begin{cases}
+1 & \text{if } p^{\mathrm{top}}_t \text{ is hit first,} \\
-1 & \text{if } p^{\mathrm{bot}}_t \text{ is hit first,} \\
\phantom{+}0 & \text{if neither is hit within } H \text{ steps.}
\end{cases}
\]

This three-class label corresponds to a real trading outcome:
``a long at this observation, held with a $\pi_u \hat\volat_t$
target and $\pi_d \hat\volat_t$ stop, would have been a winner,
loser, or timeout.'' The volatility rescaling is essential: a $2\%$
move counts differently in calm vs.\ crisis markets. Fixed-percent
barriers (``top at $+3\%$'') are trivially satisfied in high-vol
regimes and never triggered in low-vol regimes, making them
useless classification targets.

\subsection{Worked example on BTC-like daily data}
\label{subsec:ml_tb_worked}

We run the triple-barrier method on a synthetic BTC-like daily
series with ten labeled starting observations ($i = 0, \ldots, 9$),
constant volatility $\hat\volat = 2.0\%$, $\pi_u = \pi_d = 2.0$,
horizon $H = 5$ days. The full price path has sixteen entries (the
extra six are lookahead for the final starting observations):

\[
\begin{aligned}
p &= (100.00,\ 101.00,\ 102.50,\ 104.20,\ 103.80,\ 103.10,\\
&\phantom{= (}\ 102.50,\ \phantom{0}98.00,\ \phantom{0}97.50,\ \phantom{0}99.00,\ 100.20,\\
&\phantom{= (}\ 100.50,\ \phantom{0}99.80,\ 100.10,\ 100.30,\ 100.40).
\end{aligned}
\]

For each starting index $i = 0, 1, \ldots, 9$ we compute the top
barrier $p_i \cdot 1.04$, the bottom barrier $p_i \cdot 0.96$,
and walk forward up to $5$ days looking for the first barrier
hit. The resulting label table is:

\begin{center}
\begin{tabular}{c|rrr|l|c}
\toprule
$i$ & $p_i$ & top & bot & first hit & $y_i$ \\
\midrule
$0$ & $100.00$ & $104.00$ & $96.00$ & top at $j=3$, $p=104.20$ & $+1$ \\
$1$ & $101.00$ & $105.04$ & $96.96$ & timeout (no barrier) & $\phantom{+}0$ \\
$2$ & $102.50$ & $106.60$ & $98.40$ & bot at $j=5$, $p=98.00$ & $-1$ \\
$3$ & $104.20$ & $108.37$ & $100.03$ & bot at $j=4$, $p=98.00$ & $-1$ \\
$4$ & $103.80$ & $107.95$ & $99.65$ & bot at $j=3$, $p=98.00$ & $-1$ \\
$5$ & $103.10$ & $107.22$ & $98.98$ & bot at $j=2$, $p=98.00$ & $-1$ \\
$6$ & $102.50$ & $106.60$ & $98.40$ & bot at $j=1$, $p=98.00$ & $-1$ \\
$7$ & $\phantom{0}98.00$ & $101.92$ & $94.08$ & timeout (no barrier) & $\phantom{+}0$ \\
$8$ & $\phantom{0}97.50$ & $101.40$ & $93.60$ & timeout (no barrier) & $\phantom{+}0$ \\
$9$ & $\phantom{0}99.00$ & $102.96$ & $95.04$ & timeout (no barrier) & $\phantom{+}0$ \\
\bottomrule
\end{tabular}
\end{center}

The label distribution is one $+1$, five $-1$, four zero, matching
the intuitive reading: an early rally, a sharp drop triggering a
cascade of stop hits, then a late recovery where barriers are never
reached within the $5$-day window. A classifier trained on features
at each index $i$ gets a realistic, trading-aligned label, not a
raw one-day-forward return sign.

\begin{keyfactbox}[Triple-barrier method]
For each observation $t$, define a profit-target top barrier, a
stop-loss bottom barrier (both scaled by current volatility
$\hat\volat_t$), and a max-holding vertical barrier. Label the
observation with the first barrier hit: $+1$, $-1$, or $0$. The
labels match real trading outcomes and automatically adapt to
time-varying volatility. Standard reference:
\citeAFML{} Ch.~3.
\end{keyfactbox}

\begin{pitfallbox}[Don't forget to purge when using triple-barrier
labels]
Triple-barrier labels have a horizon $H$: the label at time $t$
depends on prices up to $t + H$, the exact situation that
breaks $k$-fold CV (\cref{subsec:ml_naked_return}). Every
triple-barrier classifier evaluation must use CPCV with purging
of at least $H$ observations plus an embargo. Otherwise the
reported accuracy is leakage.
\end{pitfallbox}

% =====================================================================
\section{Fractional differentiation}
\label{sec:ml_fracdiff}

The second core AFML preprocessing trick is \textbf{fractional
differentiation}. It solves a dilemma anyone training on price
series has hit.

\subsection{The dilemma}
\label{subsec:ml_fracdiff_dilemma}

Most price series are non-stationary: drifting mean, time-varying
variance. Stationarity is the assumption powering every classical
ML method: if the target distribution changes, a model fitted on
the past does not generalise. The classical fix is first differences
(log returns) $r_t = \ln p_t - \ln p_{t-1}$, which are approximately
stationary. But differencing destroys level information: returns at
$p = 100$ and $p = 10{,}000$ are indistinguishable, even though a
trader would act very differently. Return-based features have no
memory of where the price sits in its historical range, which is
exactly what a mean-reversion or range-bound strategy needs.

Fractional differentiation threads the needle: stationarity with
maximum ``memory'' preserved.

\subsection{The construction}
\label{subsec:ml_fracdiff_construction}

For a real parameter $d$, the fractional differentiation operator
is defined by the binomial expansion of $(1 - B)^d$ where $B$ is
the lag operator ($B x_t = x_{t-1}$):

\[
(1 - B)^d = \sum_{k=0}^{\infty} \binom{d}{k} (-1)^k B^k
  = \sum_{k=0}^{\infty} w_k B^k.
\]

The weights $w_k$ obey the recursion

\[
w_0 = 1, \qquad
w_k = -\,w_{k-1} \cdot \frac{d - k + 1}{k}, \quad k \ge 1.
\]

For integer $d$ this collapses to the classical differencing
operator: at $d = 1$, $w_0 = 1$, $w_1 = -1$, and all higher
weights are zero, recovering $x_t - x_{t-1}$. For fractional
$d \in (0, 1)$, the weights decay like a power law rather than
cutting off abruptly, and the operator applied to the series is
a long-memory weighted sum over the entire history of the
process:

\[
\tilde x_t = \sum_{k=0}^{\infty} w_k \cdot x_{t-k}.
\]

Lop\'ez de Prado's recipe: start at $d = 1$, decrease slowly; at
each value run an Augmented Dickey--Fuller (ADF) test on $\tilde
x_t$; the smallest $d$ for which ADF rejects non-stationarity at
5\% is optimal: just stationary enough, maximum memory. For
most equity indices this lands in $[0.3, 0.5]$, well below the
classical $d = 1$.

\subsection{Sanity-check of the weights}
\label{subsec:ml_fracdiff_numbers}

To see the long-memory behaviour concretely, here are the first
eleven weights for $d = 0.4$ (the midpoint of the typical equity
range), rounded to five decimals:

\[
w_{0..10} = (1,\ -0.4,\ -0.12,\ -0.064,\ -0.0416,\ -0.02995,
\]\[
\ -0.02296,\ -0.01837,\ -0.01516,\ -0.0128,\ -0.01101).
\]

Sum of absolute values: $\approx 1.74$. Weights decay geometrically
but slowly; the tail is long. $d = 0.3$ decays more slowly (more
memory), $d = 0.5$ faster (more stationarity). All three sit
between the extremes of ``no differencing'' ($d = 0$, max memory,
non-stationary) and ``first difference'' ($d = 1$, zero memory,
stationary).

The infinite sum is truncated: weights below $\tau \sim 10^{-5}$
are dropped, giving a finite-window implementation that loses only
a warm-up period equal to the truncation window.

\begin{keyfactbox}[Fractional differentiation]
For $d \in (0, 1)$, the fractionally differenced series
$\tilde x_t = \sum_{k=0}^{\infty} w_k x_{t-k}$ with
$w_0 = 1$, $w_k = -w_{k-1}(d - k + 1)/k$ achieves stationarity
with minimum loss of memory. The optimal $d$ is the smallest
value for which the series passes an ADF test; for most equity
return data it falls in $[0.3, 0.5]$. Standard reference:
\citeAFML{} Ch.~5.
\end{keyfactbox}

% =====================================================================
\section{Information-driven bars}
\label{sec:ml_bars}

A hidden assumption of every textbook time-series ML tutorial is
uniform time sampling: one-minute, one-day, or one-hour bars. This
is the wrong scheme for trading data. Information does not arrive
uniformly: most of the day is quiet, then a news release dumps
five minutes' information into five seconds. Uniform-time sampling
means irregular sampling on the news clock: over-sampling quiet
periods, under-sampling active ones.

\citeAFML{} Ch.~2 proposes activity-adaptive sampling. Four
standard families, in order of sophistication:

\subsection{Tick bars}
\label{subsec:ml_ticks}

A \textbf{tick bar} closes every $N$ trades. Each bar represents
the same transaction count regardless of elapsed time. Quiet
periods produce fewer bars, news spikes many. The resulting series
has far more stable intraday return distributions than time bars,
which is exactly what classical ML assumptions want.

\subsection{Volume bars}
\label{subsec:ml_volume}

A \textbf{volume bar} closes every $V$ shares (or contracts).
Trades of $100$ and $900$ shares trigger a close at $V = 1000$;
a single $2000$-share trade triggers two closes in succession.
Volume bars respect that large trades carry more information than
small ones; tick bars ignore this.

\subsection{Dollar bars}
\label{subsec:ml_dollar}

A \textbf{dollar bar} closes every $D$ dollars of notional. If
the price doubles, half as many shares represent the same
dollar amount, so the bar boundary adapts. For equities that
appreciate or depreciate substantially, dollar bars give the
most stable return statistics across time, more so than
volume bars (which drift with share-price inflation) or tick
bars.

\subsection{Information-driven bars}
\label{subsec:ml_info_bars}

\textbf{Information-driven bars} go further: a bar closes when a
cumulative information statistic exceeds a threshold. The
simplest example is the signed-volume imbalance bar: accumulate
$\sum (\text{sign}(r_t) \cdot v_t)$ until $|\cdot| > \theta$. The
bar closes when flow becomes directionally significant. Balanced
chop produces few bars; a large directional order produces many
in rapid succession. Switching from time to dollar (or
information-driven) bars is often the most effective single change in
a feature-engineering pipeline: the same classifier on the same
features improves because the time axis is now aligned with the
clock on which information actually arrives.

\begin{keyfactbox}[Bar taxonomy]
\begin{itemize}
  \item \textbf{Time bars}: fixed wall-clock interval. Worst
  for ML; the default in every naive tutorial.
  \item \textbf{Tick bars}: fixed number of trades. Better;
  equalises transaction count.
  \item \textbf{Volume bars}: fixed shares/contracts. Better
  still; weights by trade size.
  \item \textbf{Dollar bars}: fixed notional. Best plain
  alternative; robust to price drift over time.
  \item \textbf{Information-driven bars}: variable, adapts to
  cumulative signed flow or other entropy proxy. Sampling
  density tracks information arrival.
\end{itemize}
Any ML-for-trading pipeline that still uses one-minute bars
after reading \citeAFML{} Ch.~2 is leaving a free lunch on the
table.
\end{keyfactbox}

% =====================================================================
\section{Meta-labeling: the single best ML-for-trading idea}
\label{sec:ml_metalabel}

If you take one idea from this chapter, take this one.
Meta-labeling is cleanly positive-expected-value for any
researcher with a working rule-based primary strategy. The
construction is embarrassingly simple, which is partly why the
field was slow to notice it.

\subsection{Setup}
\label{subsec:ml_meta_setup}

Suppose you have a \textbf{primary strategy} emitting trade
signals, e.g.\ the OU mean-reversion rule from
\cref{ch:mean_reversion_ou}: long on $z < -2$, short on $z > +2$.
Suppose its empirical hit rate is $55\%$, annualised Sharpe $1.0$,
drawdowns occasionally reaching $20\%$: decent but not great,
painful in specific regimes.

The meta-labeling question: \emph{can we learn to skip the bad
triggers?} Instead of predicting returns, we predict whether a
primary signal, given current market state, is a winner or loser.
We act on signals the classifier flags ``probably winner'' and
skip the rest.

\subsection{The mechanics}
\label{subsec:ml_meta_mechanics}

The meta-labeling pipeline has four steps.

\begin{enumerate}
  \item \textbf{Run the primary strategy on historical data}
  and log every trigger: timestamp, direction (long/short),
  and eventual outcome (``winner: closed in profit before
  stop'' or ``loser: closed in loss''). This gives a binary
  label $m_t \in \{0, 1\}$ for every primary trigger.
  \item \textbf{Build features for each trigger} from the
  current market state: regime (low-vol / high-vol /
  crisis), recent volatility, top-of-book imbalance, time
  of day, spread, and so on. These are the same kind of
  features as in \cref{subsec:ml_feature_eng}.
  \item \textbf{Train a binary classifier}
  $\hat m(\mathbf{x})$ (logistic regression, random forest,
  gradient boosting; the choice of model matters less
  than the features) to predict $m_t$ from the features
  $\mathbf{x}_t$, \emph{under CPCV}.
  \item \textbf{Deploy}: at each new primary trigger, compute
  $\hat m(\mathbf{x}_t)$; if the predicted probability is
  above a threshold $\tau$ (say $0.55$), take the trade;
  otherwise pass.
\end{enumerate}

The classifier does not predict returns. It only improves on the
$55\%$ base rate by preferentially skipping losers. This is a
far easier problem: clean binary target, reasonable class
balance, features natively predictive of trade quality
(volatility regime, liquidity, time of day). Above all, the
learner is not asked to generate alpha. It is asked to
\emph{gate} alpha that already exists.

\subsection{Worked lift calculation}
\label{subsec:ml_meta_lift}

Consider a stylised pairs-trading strategy
(\cref{ch:pairs_cointegration}) generating $N = 1000$ primary
triggers over a historical window. Each winner contributes $+1R$,
each loser $-1R$ ($R$ = per-trade risk unit):

\[
\text{Baseline P\&L} = (0.55 \cdot 1000) \cdot (+1R)
  + (0.45 \cdot 1000) \cdot (-1R) = 100R.
\]

Train a meta-labeler that flags $40\%$ of trades ($400$ triggers)
as skip. Assume it is informative but not perfect: of the $400$
skipped, $240$ were losers, $160$ were winners (vs.\ a
uniform-skip baseline of $220$ losers, $180$ winners: a modest
edge). We keep $600$ trades, of which

\[
\text{kept winners} = 550 - 160 = 390, \qquad
\text{kept losers} = 450 - 240 = 210.
\]

The kept hit rate is $390 / 600 = 65\%$, a $10$-pp lift from
$55\%$. Kept P\&L:

\[
(390)(+1R) + (210)(-1R) = 180R,
\]

\emph{an $80\%$ lift in absolute P\&L}, on $40\%$ fewer trades.
Assuming per-trade variance is unchanged and the
$\sqrt{N}$ denominator rescales with trade count, the stylised
Sharpe jumps from $100 / \sqrt{1000} \approx 3.16$ to
$180 / \sqrt{600} \approx 7.35$, a factor of $\sim 2.3\times$.
These numbers are illustrative, but they show why meta-labeling
earns its billing: a modest classifier merely
correlated with trade quality translates into dramatic
bottom-line improvement, because all it does is filter.

\subsection{Why meta-labeling works (and return prediction doesn't)}
\label{subsec:ml_meta_why}

Meta-labeling works where naked return prediction fails because
it attacks a fundamentally easier problem. Return prediction asks
which of all market states will have positive next-step returns;
the signal-to-noise is terrible and the target is a
near-random walk. Meta-labeling asks, given that the primary just
fired, which firings are the good ones. This conditioning is the
trick: the primary has already isolated high-information moments,
and the classifier only has to rank them. The target is binary and
stable (``did this entry hit target before stop?''), the feature
set is tight, class balance is sane, and the learner has no
ambition to beat the market from scratch, only to improve a
working baseline.

\begin{keyfactbox}[Meta-labeling]
Given a rule-based primary strategy with a modest positive
edge, a binary classifier trained on trade-quality features
can systematically skip the worst triggers and keep the best,
turning a mediocre strategy into a good one without touching
the primary logic. This is the single most valuable idea
in ML for trading. Standard reference: \citeAFML{} Ch.~3.
\end{keyfactbox}

\begin{intuitionbox}
Meta-labeling inverts the usual framing. Instead of ``predict
returns, then construct a strategy,'' it says ``construct a
strategy, then learn to filter its triggers.'' The strategy
generates alpha; the classifier chooses battles. It works
because the primary strategy has collapsed the
infinite-dimensional ``when is the market predictable'' down to
a finite set of concrete trade opportunities, each with a
binary supervisable outcome. That is a far better shaped ML
problem.
\end{intuitionbox}

\begin{pitfallbox}[Don't meta-label until the primary strategy is
profitable]
Meta-labeling can only filter; it cannot turn a losing primary
into a winner. If the primary has a $45\%$ hit rate, the best
meta-labeling can do is skip every trade (breakeven via not
trading). The literature sometimes sells it as magic, but the
magic is conditional on the primary having an edge. Validate the
primary with CPCV (\cref{ch:backtesting}) first; only then build
a meta-labeler.
\end{pitfallbox}

% =====================================================================
\section{The hardest problem is not prediction, it is validation}
\label{sec:ml_validation_is_hard}

Standard ML research orders things so: prediction is hard (train a
model with low training error), validation is the easy final step
(held-out set, report accuracy). For trading, that ordering is
inverted. The hard part is \emph{validation}: under which
measurement procedure can you trust that reported performance
survives contact with real money? Fitting the model is a footnote.

Why? A trading ML model has three adversaries a normal one does
not: (i)~the target is weakly predictable, so the honest OOS
signal is drowned by any validation bug; (ii)~the data is
non-stationary, so one held-out period is an unreliable estimate;
(iii)~a broken model loses money, unrecoverable by ``a second
experiment with cleaner methodology.''

Consequence: the majority of time goes to the validation pipeline,
not the model. Every technique in \cref{ch:backtesting}
(walk-forward, CPCV with purging and embargo, deflated Sharpe,
nested CV, plateau-not-peak selection) is required by default.
Every technique here (triple-barrier labels, fractional
differentiation, information-driven bars, meta-labeling) sets
up the problem so validation has a fighting chance.

\begin{keyfactbox}[The validation-first doctrine]
In ML for trading the ranking of effort is:
\begin{enumerate}
  \item Validation pipeline (CPCV, purging, embargo,
  nested CV, deflated Sharpe).
  \item Feature engineering (order-book microstructure,
  volatility normalisation, information-driven bars).
  \item Label construction (triple-barrier, meta-labeling
  targets).
  \item The model itself (and this step is nearly always
  ``a gradient-boosted tree with sensible defaults'').
\end{enumerate}
The academic ranking is the reverse. If you are spending most
of your time on step~4, you are working on the wrong problem.
\end{keyfactbox}

% =====================================================================
\section{Prosperity: ML plays a narrow role}
\label{sec:ml_prosperity}

\begin{prosperitybox}
In IMC Prosperity, the top teams' edge is almost never ML. Profit
comes from classical tools: market making
(\cref{ch:market_making}), mean reversion and pairs trading
(\cref{ch:mean_reversion_ou,ch:pairs_cointegration}), basket/ETF
arbitrage (\cref{ch:basket_etf_arb}), and options
(\cref{ch:black_scholes,ch:greeks_hedging}). Neural networks on
tick features do not appear in winning writeups. The reason is
specific: Prosperity data is short (a few hundred thousand
observations per round at most), products are synthetic with
stationary microstructure, and classical-strategy edge is always
large enough to statistically dominate an ML model fitted on a
few thousand ticks.

One area where ML-flavoured technique does appear:
\textbf{detecting informed traders from per-trade data}. The
Frankfurt Hedgehogs' second-place Prosperity~3 writeup (source of
\cref{sec:olivia}'s Olivia story) trained a pattern-detector on
per-trader behaviour to identify anonymous bots filling at daily
extremes. In this chapter's framing: a classification problem.
Given history labeled by anonymous trader ID, predict ``is this
trader informed?'' Features are classical: fill-price distance
from the day's running min/max, fill size, time since last trade.
The binary label derives from the observation that one bot always
fills at the daily extreme. The classifier is used as a
\emph{meta-labeler} on the team's own trades: ``don't compete
with informed flow, align with it.''

This is precisely the meta-labeling idea of
\cref{sec:ml_metalabel}. The primary is a rule-based pairs trade
on an ETF basket; the meta-labeler biases entry thresholds by the
inferred informed-bot direction. The reported lift materially
improved placement, and the classical part still worked when the
informed bot did not trade; ML was pure upside, not a
replacement.

General lesson: do not reach for ML until classical tools are
exhausted. When you do, reach for a one-sentence classification
problem (``is this trader informed,'' ``is this regime
high-volatility,'' ``is this spread about to revert''), not
``predict the next tick.'' Prosperity punishes overfitting fast,
and the surviving techniques are the boring ones: few features,
linear or tree-based models, CPCV.
\end{prosperitybox}

% =====================================================================
\section{Goldman Sachs: which desks actually use ML}
\label{sec:ml_gs}

\begin{gsbox}
At Goldman Sachs, as at any large sell-side bank, ML lives on
specific desks and problems; it is not the bank-wide religion
outsiders sometimes imagine. The distribution maps cleanly onto
the ``where ML helps'' list in \cref{sec:ml_where_helps}:

\begin{itemize}
  \item \textbf{Systematic Equities (quantitative
  long/short)}: the cleanest home for meta-labeling.
  Quant research generates primary alphas from
  fundamental or technical signals; the systematic
  trading team trains classifiers on trade-level
  features to filter the worst entries. Alphas are
  transparent and attributable; ML is a conditional
  filter on top.
  \item \textbf{Central Risk Book}: the team
  warehousing inventory across flow desks uses HMMs and
  related regime-detection models to time hedging bands
  (high-vol regime, widen; quiet regime, tighten). ML
  here is a regime classifier, not a return predictor.
  \item \textbf{Execution / Algorithmic Trading}: uses
  RL for child-order placement atop an Almgren--Chriss
  baseline (\cref{ch:almgren_chriss}), plus supervised
  learning for short-horizon fill-probability prediction
  in smart-order routers. The highest-compute ML use
  case in the firm.
  \item \textbf{Data Science / Strats}: builds
  feature-engineering pipelines on alternative and
  microstructure data (credit-card transactions,
  satellite imagery, news sentiment, order-book
  reconstructions), sold as \emph{inputs} to the
  systematic desks. Output is ``a feature matrix,'' not
  ``a model.''
\end{itemize}

Three points worth underlining. First, \textbf{ML is almost
never used for directional signal generation}. The directional
signal, long or short, comes from quantitative research
(QR), using classical statistical and economic models with
defensible attributions. ML enters as filter, regime switch, or
execution optimiser, never to generate the thesis. A Strat
pitching an RNN return predictor as a trading signal is gently
but firmly redirected to an HMM or linear factor model. Second,
\textbf{validation discipline is taken very seriously}: every
production ML model has a CPCV evaluation, a monthly-recut
out-of-sample window, and a documented retirement criterion
(e.g., ``pull if rolling 6-month OOS Sharpe $< 0.5$''). Models
do not live forever. Third, \textbf{the feature matrix is the
crown jewel}: the most durable source of desk edge is a
well-engineered feature pipeline that can be recombined as
regimes change. Fancy model on a bad matrix loses to a simple
model on a good one, every time.
\end{gsbox}

% =====================================================================
\section{A short reference implementation}
\label{sec:ml_reference_impl}

One short Python snippet for the triple-barrier method is this
chapter's core algorithmic contribution; everything else is about
what to feed in and how to validate what comes out. The
implementation below takes a price array, per-observation
volatility, profit/stop multipliers, and max holding period, and
returns the triple-barrier label per observation.

\begin{lstlisting}[language=Python, caption={Triple-barrier
labeling, producing the label table of
\cref{subsec:ml_tb_worked}.}]
import numpy as np

def triple_barrier_labels(prices, sigma, pt_mult=2.0, sl_mult=2.0, max_hold=5):
    """Return +1 (top), -1 (stop), 0 (timeout) for each observation."""
    n = len(prices)
    labels = np.zeros(n)
    for i in range(n - max_hold):
        top = prices[i] * (1 + pt_mult * sigma[i])
        bot = prices[i] * (1 - sl_mult * sigma[i])
        for j in range(1, max_hold + 1):
            if prices[i + j] >= top:
                labels[i] = +1; break
            if prices[i + j] <= bot:
                labels[i] = -1; break
    return labels
\end{lstlisting}

Run this on the ten-observation BTC-like series of
\cref{subsec:ml_tb_worked} with constant \texttt{sigma[i] = 0.02}
and you recover the label vector
$(+1, 0, -1, -1, -1, -1, -1, 0, 0, 0)$, exactly the worked
example. An implementation reproducing this vector has the
barrier logic right; one that does not has a sign error, an
off-by-one, or a confused volatility-scaling direction. Production versions (e.g.,
\texttt{mlfinlab.labeling.get\_events} in the \texttt{mlfinlab}
package accompanying \citeAFML{}) add volatility estimation,
asymmetric horizons, and event metadata, but the core logic is
the ten lines above.

% =====================================================================
\section{Key facts to memorise}
\label{sec:ml_key_facts}

\begin{itemize}
  \item \textbf{ML for trading is narrow.} It helps with
  feature engineering on order-book data, regime
  detection, trade-level meta-labeling, and execution. It
  does not help with naked return prediction from lagged
  returns, no matter how deep the network.
  \item \textbf{Naive $k$-fold CV on time series is a bug.}
  Triple-barrier labels with horizon $H$ leak between
  folds unless you purge training observations whose
  label windows intersect test observations (plus an
  embargo). Use CPCV from \cref{ch:backtesting}.
  Everything else in this chapter assumes this fix is
  already applied.
  \item \textbf{Triple-barrier labels} define profit,
  stop, and timeout barriers for each observation and
  label it by the first barrier hit. Volatility scaling
  is mandatory; fixed percentage barriers are useless.
  Standard reference: \citeAFML{} Ch.~3.
  \item \textbf{Fractional differentiation} threads the
  needle between ``stationary but memoryless'' returns
  and ``non-stationary but level-preserving'' prices.
  Weights: $w_0 = 1$,
  $w_k = -w_{k-1}(d - k + 1)/k$. Typical optimal $d$
  for equity series is in $[0.3, 0.5]$, chosen as the
  smallest $d$ at which an ADF test rejects
  non-stationarity.
  \item \textbf{Information-driven bars} (tick, volume,
  dollar, imbalance) sample time more densely when
  something is happening. Switching from time bars to
  dollar bars is often the most effective single
  change in a feature pipeline.
  \item \textbf{Meta-labeling is the single
  most valuable idea.} Train a binary classifier on
  top of a working rule-based primary strategy to filter
  out low-quality triggers. Turns a mediocre strategy
  into a good one without touching the primary logic.
  \item \textbf{Meta-labeling only filters.} It cannot
  turn a losing primary strategy into a winner. Validate
  the primary with CPCV first.
  \item \textbf{The validation pipeline is the hard
  part.} In ML for trading the ranking of effort is
  (1)~validation, (2)~features, (3)~labels, (4)~model.
  The academic ranking is the reverse.
  \item \textbf{Non-stationarity never goes away.}
  Rolling retraining, regime conditioning, and deflated
  Sharpe all soften it; none eliminate it. Any model
  expected to live for years is the wrong model.
  \item \textbf{ML at Goldman Sachs} lives on systematic
  equities (meta-labeling), central risk
  (regime detection), execution (RL for child orders),
  and Strats data science (feature pipelines). It does
  \emph{not} live in the directional signal: that comes
  from QR using classical models.
  \item \textbf{ML in Prosperity is rarely the edge.}
  Classical tools dominate; the one place ML appears is
  in \emph{informed-trader detection} (Olivia-style),
  which is a meta-labeling application in disguise.
  \item \textbf{The single sentence.} \emph{``The hardest
  problem in ML for trading is not prediction, it is
  validation.''} Everything in this chapter is in
  service of that sentence.
\end{itemize}
